% Construcción dodecafásica del funcional de retorno.
\section{Cadena fundamental dodecafásica y funcional hexádico--octádico}

La incidencia activa se construye antes de toda realización gravitatoria.
Sea
\[
 H_{\mathrm{act}}=\{1,2,4,5,7,8\},
 \qquad
 \mathcal P_2(H_{\mathrm{act}})
 =\{A\subset H_{\mathrm{act}}:|A|=2\}.
\]
Entonces
\[
 |H_{\mathrm{act}}|=6,
 \qquad
 |\mathcal P_2(H_{\mathrm{act}})|=15,
\]
y las incidencias anterior y posterior a la ampliación poseen cardinales
\[
 \mathcal F_{90}=H_{\mathrm{act}}\times
 \mathcal P_2(H_{\mathrm{act}}),
 \qquad |\mathcal F_{90}|=90,
\]
\[
 \mathcal F_{120}=\{1,\ldots,8\}\times
 \mathcal P_2(H_{\mathrm{act}}),
 \qquad |\mathcal F_{120}|=120.
\]
Estos dos cardinales son salidas de la regla de fase; no son ángulos ni
parámetros recibidos desde la realización física.

Denotemos por \(C_{12}=\mathbb Z/12\mathbb Z\) el registro temporal del TPK y
por \(\partial_{\rm ret}\) su única costura de retorno por vuelta. En el
módulo libre de sucesos orientados se considera la cadena fundamental
\begin{equation}
 c_{12}^{+}
 =\sum_{j\in C_{12}}[j,\mathcal F_{90}]
  -[\partial_{\rm ret},\mathcal F_{120}].
 \label{intsuccpass:eq:c12-fundamental-chain}
\end{equation}
La suma contiene exactamente los doce sectores de la rueda; el segundo
término es único porque una vuelta posee una sola costura de cierre. Su signo
es el signo de frontera. La involución bidireccional envía
\(c_{12}^{+}\) a \(c_{12}^{-}=-c_{12}^{+}\): cambia la orientación de la
cadena, pero no altera sus multiplicidades absolutas ni sus tipos de
incidencia.

Para \(m\in\{90,120\}\), sea
\[
 S_m(q)=\frac{q^m}{1-q^{3m}}.
\]
El lector de retorno \(\mathscr R\) se define sobre los generadores de
\eqref{intsuccpass:eq:c12-fundamental-chain} por
\[
 \mathscr R([j,\mathcal F_{90}])=S_{90},
 \qquad
 \mathscr R([\partial_{\rm ret},\mathcal F_{120}])=S_{120}.
\]
La primera igualdad es constante sobre \(C_{12}\) por covariancia bajo la
rotación de fase; la segunda tiene un dominio distinto y sólo actúa en la
costura ampliada. Por linealidad,
\begin{equation}
 \boxed{\mathscr R(c_{12}^{+})=12S_{90}-S_{120}.}
 \label{intsuccpass:eq:dodecaphase-forces-12minus1}
\end{equation}
Así, los coeficientes \(12:-1\) son la imagen de la cadena fundamental: doce
celdas de fase y una frontera orientada. No se seleccionan mediante el valor
convencional del parámetro de Barbero--Immirzi, ni mediante una comparación
decimal, ni imponiendo previamente el coeficiente de polo que se calculará a
continuación.

\begin{theorem}[Coeficiente de polo de la cadena dodecafásica]
El funcional \(F(q)=\mathscr R(c_{12}^{+})(q)
=12S_{90}(q)-S_{120}(q)\), determinado por los doce sectores
y su costura orientada, tiene coeficiente \(1/24\) en
\((1-q)^{-1}\). Su residuo complejo en \(q=1\) es \(-1/24\).
El coeficiente de polo es una salida del funcional ya construido.
\end{theorem}

\begin{proof}
Para un funcional con polo simple en \(q=1\), escribimos
\[
 p_1(F):=\lim_{q\uparrow1}(1-q)F(q).
\]
La identidad
\[
 1-q^{3m}=(1-q)\sum_{j=0}^{3m-1}q^j
\]
implica
\[
 p_1(S_m)
 =\lim_{q\uparrow1}\frac{q^m}{\sum_{j=0}^{3m-1}q^j}
 =\frac1{3m}.
\]
Por linealidad,
\[
 p_1\bigl(\mathscr R(c_{12}^{+})\bigr)
 =\frac{12}{270}-\frac1{360}
 =\frac{48-3}{1080}=\frac1{24}.
\]
Estas igualdades se efectúan en \(\mathbb Q\). El residuo complejo
es el coeficiente de \((q-1)^{-1}\), de modo que
\[
 \operatorname*{Res}_{q=1}\mathscr R(c_{12}^{+})(q)
 =-p_1\bigl(\mathscr R(c_{12}^{+})\bigr)=-\frac1{24}.
\]
\end{proof}

La ecuación diofántica
\[
 \frac a{270}-\frac b{360}=\frac1{24}
 \quad\Longleftrightarrow\quad
 4a-3b=45
\]
permanece como certificado aritmético posterior. Por sí sola admite la
familia \((a,b)=(12+3t,1+4t)\); la cadena
\eqref{intsuccpass:eq:c12-fundamental-chain} elimina esa ambigüedad antes del cálculo,
pues cambiar \((12,1)\) altera la multiplicidad de la órbita o de su costura
y, por tanto, ya no representa una vuelta fundamental del TPK.

Con los canales conjugados \(q_\pm\) ya generados por la rama angular, se
obtiene finalmente
\[
 \mathscr R(c_{12}^{+})(q_-)-
 \mathscr R(c_{12}^{+})(q_+)
 =12\Delta S_{90}-\Delta S_{120}.
\]
Esta es la entrada interna del funcional hexádico--octádico. Su evaluación
como parámetro de Barbero--Immirzi y su inserción en el operador de área son
publicaciones posteriores: no participan en la selección de
\(c_{12}^{+}\), de \(90/120\), de \(12:-1\) ni de \(1/24\).

\paragraph{Falsadores.}
La construcción queda refutada si el censo de fase no da doce sectores, si
una vuelta contiene más de una costura, si las incidencias no tienen
cardinales \(90\) y \(120\), si la involución no invierte la cadena completa,
o si el coeficiente exacto de \((1-q)^{-1}\) de su imagen difiere de \(1/24\).
